\documentclass[11pt,a4paper]{article}
\usepackage[T1]{fontenc}
\usepackage{lmodern,amsmath,amssymb,amsthm,booktabs}
\usepackage[margin=26mm]{geometry}
\usepackage[hidelinks]{hyperref}
\usepackage{xurl}
\usepackage{microtype}
\newtheorem{theorem}{Theorem}
\newtheorem{lemma}{Lemma}
\title{The whole-prefix recurrence in Conjecture 00000000209}
\author{zbw123456}
\date{13 October 2026}
\begin{document}
\maketitle
\begin{abstract}
The recurrence printed in both language versions of Conjecture 00000000209
maximizes the exponent in a representation of the entire preceding word as
$Y^k$. Its unique solution is $1,1,2,1,1,1,\ldots$; hence 5 never occurs.
We prove existence, uniqueness, and the exact formula, and supply a
self-contained, compiled Lean 4 verification.
\end{abstract}

\section{Scope and the definition discrepancy}
The checked problem is at the following official repository commit~\cite{problem}:
\begin{center}\small\texttt{1e3f7c25616e84b611350a412febceabad75b56a}\end{center}
Its displayed repetition rule is $Y^k$, applied to the \emph{entire}
prefix. No preliminary word $X$ appears in either language version.

This differs from the standard Gijswijt recurrence, which maximizes $k$
in $XY^k$, where $X$ may be empty and $Y$ is nonempty~\cite{gijswijt}.
The standard sequence is unbounded. The counterexample below concerns only
the rule actually displayed in the competition, and does not refute that
classical result. The parenthetical occurrence information in the problem
refers to the standard sequence and conflicts with the displayed rule.
\textbf{If the intended rule is corrected to $XY^k$, this submission does
not resolve the corrected problem.} This semantic issue remains for the
maintainer to assess; no official acceptance is claimed.

\section{Precise statement}
Use zero-based indexing in this section and in Lean. For $a:\mathbb N\to
\mathbb N$ and $N,k\geq1$, write $P(a,N,k)$ if there is $m\geq1$ such that
\[
 N=mk,\qquad
 a(jm+i)=a(i)\quad(0\leq j<k,\ 0\leq i<m).
\]
These are exactly the conditions that the word of length $N$ is $k$
copies of its first block of length $m$. All displayed indices are below
$N$, since $jm+i<(j+1)m\leq km=N$.

The literal recurrence is
\[
 a(0)=1,\qquad
 a(N)=\max\{k\geq1:P(a,N,k)\}\quad(N\geq1).                 \tag{R}
\]
The maximum exists: $k=1$ is admissible, and $P(a,N,k)$ implies $k\leq N$.
The Lean predicate \texttt{IsNext} expresses both attainment and maximality.
The recurrence has natural-valued outputs because every exponent is a
positive integer; the choice of this codomain adds no restriction to (R).

\begin{theorem}
There is exactly one sequence satisfying (R). It is
\[
 b(N)=\begin{cases}2,&N=2,\\1,&N\ne2.\end{cases}
\]
In particular, no sequence satisfying (R) contains 5.
\end{theorem}

\begin{proof}
First check that $b$ satisfies the recurrence. The initial value is 1.
The word of length 1 is $(1)$, so its largest repetition exponent is 1.
The word of length 2 is $(1,1)$, so its largest exponent is 2.
For every $N\geq3$, the prefix contains exactly one occurrence of 2.
If it equals $Y^k$, the number of occurrences of 2 is
$k$ times the number of occurrences of 2 in $Y$. Thus $k=1$.
Exponent 1 is always attained, proving every recurrence equation for $b$.

For clarity, Lean checks the last argument directly on indices.
Assume $k\geq2$. If $m=1$, then $N\geq3$ implies $k\geq3$, and the
third copy would give $b(2)=b(0)$, a contradiction.
If $m=2$, the second copy similarly gives $b(2)=b(0)$.
If $m\geq3$, the second copy instead gives $b(m+2)=b(2)$, also impossible.
These cases exhaust all positive block lengths.

To prove uniqueness, let $a$ and $c$ satisfy (R), and use strong induction
on $N$. Their zeroth terms are equal. If their terms below $N$ agree,
their length-$N$ prefixes admit exactly the same repetition exponents.
Both $a(N)$ and $c(N)$ are attained maxima of this same set, so each is
at most the other. Hence $a(N)=c(N)$. Taking $c=b$ proves the exact formula.
Its only values are 1 and 2, so 5 is absent.
\end{proof}

\section{Formal verification and reproducibility}
The project uses Lean \texttt{v4.19.0}, imports only \texttt{Init}, and has
no third-party dependencies. In the \texttt{lean} directory run:
\begin{verbatim}
lake build
lake env lean -DwarningAsError=true Conjecture209.lean
\end{verbatim}
The second command also prints the axiom dependencies of the main results.
The source contains no proof placeholders, no native decision procedure,
and no custom axiom declarations. The printed dependency lists contain
only the standard Lean foundations \texttt{propext} and \texttt{Quot.sound}.

\begin{center}
\small
\begin{tabular}{@{}lp{7.5cm}@{}}
\toprule
Lean result & Mathematical content\\
\midrule
\texttt{candidate\_follows\_rule} & All recurrence equations hold for $b$.\\
\texttt{rule\_unique} & Any two solutions agree at every index.\\
\texttt{exact\_formula} & Every solution has the displayed formula.\\
\texttt{no\_rule\_solution\_contains\_five} & There is no solution together with an index whose value is 5.\\
\texttt{counterexample} & A solution exists and has no occurrence of 5.\\
\bottomrule
\end{tabular}
\end{center}
The final result is universal in the index. It is not inferred from a
finite computation, nor is a hypothetical recurrence solution assumed to
exist without construction. Verification logs accompany this report.

\begin{thebibliography}{9}
\bibitem{problem}
The Last Math Competition, Conjecture 00000000209, official repository,
commit specified above.
\url{https://github.com/The-Last-Math-Competition/The-Last-Math-Competition/blob/1e3f7c25616e84b611350a412febceabad75b56a/conjectures/00000000209.md}.
\bibitem{gijswijt}
F. J. van de Bult, D. C. Gijswijt, J. P. Linderman, N. J. A. Sloane,
and A. R. Wilks, \emph{A Slow-Growing Sequence Defined by an Unusual
Recurrence}, Journal of Integer Sequences 10 (2007), Article 07.1.2.
\url{https://cs.uwaterloo.ca/journals/JIS/VOL10/Sloane/sloane55.html}.
\end{thebibliography}
\end{document}
